\documentclass[conference]{IEEEtran}
\usepackage{amsmath}
\usepackage{amssymb}
\usepackage{cite}
\usepackage{hyperref}

\begin{document}

\title{Data-Driven Priors and Physics-Based Corrections: A Hybrid Flow
Matching Approach to Safe Robot Motion Planning}

% No need to put student ID here
\author{
  \IEEEauthorblockN{Zihao Sun}
  \and
  \IEEEauthorblockN{Charles Beck}
}

\maketitle

\begin{abstract}
  UCLA ECE M237 Spring 2026 Course Project Milestone Reports. Week
  3---Modeling \& Sketch of Method.
\end{abstract}

% -------------------------------------------------------
\section{Project Title}

% Provide your project title here.

\noindent Data-Driven Priors and Physics-Based Corrections: A Hybrid
Flow Matching Approach to Safe Robot Motion Planning.

% -------------------------------------------------------
\section{Main Idea}

% Explain the main idea of your project in \textit{one sentence}. This
% can be copied from your previous reports, but also feel free to update
% it if you have refined your idea.

We investigate how to merge an \textbf{offline}-trained flow matching
trajectory prior (learned from demonstration data) with an
\textbf{online} model-based score correction (computed analytically
from the known dynamics and cost), and study whether the resulting
hybrid yields trajectories that are simultaneously diverse and
physically feasible.

% -------------------------------------------------------
\section{Problem Formulation}

% Describe in \textit{one paragraph} the problem you want to solve. I
% encourage to use equations to describe your problem.

Let $\tau=(x_0,u_0,\ldots,x_{H-1},u_{H-1},x_H)$ be a trajectory of
the discrete-time system $x_{k+1}=f(x_k,u_k)$, with running cost
$J(\tau)=\sum_{k=0}^{H-1}\ell(x_k,u_k)+\ell_H(x_H)$. Given a
demonstration dataset
$\mathcal{D}=\{\tau^{(i)}\}_{i=1}^{N}\!\sim\!p^{\mathrm{data}}$, our
goal is to sample from the product-of-experts target
$p^{\mathrm{target}}(\tau)\propto p^{\mathrm{data}}(\tau)\,
[p^{\mathrm{model}}(\tau)]^{\lambda}$, where
$p^{\mathrm{model}}(\tau)\propto\exp(-J(\tau)/\beta)$ is the
Boltzmann distribution induced by the known dynamics and cost. At
each diffusion time $t$ this is equivalent to denoising with the
composite score
\begin{equation}
  \nabla\log p^{\mathrm{target}}_t \approx
  \nabla\log p^{\mathrm{data}}_t
  + \lambda(x,t)\,\nabla\log p^{\mathrm{model}}_t,
  \label{eq:score}
\end{equation}
which is the central object our project formalizes and analyzes; a
CBF safety term $\mu\,\nabla\log p^{\mathrm{safe}}_t$ is an optional
extension.

% -------------------------------------------------------
\section{Sketch of Method}

\subsection{Offline Component}

% Describe in \textit{one paragraph} the offline component of your
% method (for instance, which computes the value function or the
% policy). Explain why it is suitable for your problem. Discuss the
% computation complexity.

Offline, we train a flow matching velocity field $v_\theta(\tau_t,t)$
on $\mathcal{D}$ via the conditional flow matching
loss~\cite{lipman2023flow}
$\mathcal{L}_{\mathrm{CFM}}=\mathbb{E}_{t,\tau_0,\tau_1}\|
v_\theta((1-t)\tau_0+t\tau_1,t)-(\tau_1-\tau_0)\|^2$, with $\tau_t$
linearly interpolating Gaussian noise and a demonstration. The
flow--score identity~\cite{song2021score} then yields an implicit
data-driven score
$s_\theta\!\approx\!\nabla\log p^{\mathrm{data}}_t$ at inference, in
line with diffusion-based planners~\cite{flowmp,janner2022}.
Amortizing the multimodal behavior of demonstrations into a network
replaces $O(MH)$ rollouts with one $O(|\theta|)$ forward pass per
step---an offline analogue of fitted value iteration, paid once at
$O(N\!\cdot\!\text{epochs}\!\cdot\!|\theta|)$.

\subsection{Online Component}

% Describe in \textit{one paragraph} the online component of your
% method, which will be used for the actual policy execution. Explain
% why it is suitable for your problem. Discuss why it is necessary or
% could be helpful in complementing the limitations of the offline
% part.

At inference we evaluate a model-based correction
$\nabla\log p^{\mathrm{model}}_t\!\approx\!
\mathbb{E}_{w(\tau)}[\tau-\tau_t]/\sigma_t^2$ from self-normalized
importance sampling of cost-weighted rollouts
$w(\tau)\propto\exp(-J(\tau)/\beta)$, in the spirit of Model-Based
Diffusion~\cite{mbd}, and add it to $s_\theta$ at every step. This
is essential because the offline prior is blind to the deployed cost
and to perturbations of dynamics or environment, and degrades out of
distribution. MBD itself suffers from high estimator variance over
long horizons and lacks a data prior; we mitigate both by anchoring
rollouts at the flow-matching posterior mean and letting
$\lambda(x,t)$ react to local data density via a Gaussian fit to
U-Net features. In the LQ case the correction reduces to
$-P_t x_t$, an LQR-style sanity check we plan to raise with the MBD
authors at the guest lecture.

% -------------------------------------------------------
\section{Working Example Modeling}

\subsection{Modeling \& Simulation}

% Describe in \textit{one paragraph} the \textit{model} of the running
% example system on which you will deploy your method.

We instantiate the framework on a 7-DoF Franka Panda doing tabletop
reaching in \texttt{MuJoCo}: state
$x=(q,\dot q)\!\in\!\mathbb{R}^{14}$, torque
$u\!\in\!\mathbb{R}^{7}$ saturated at $\pm 87$\,Nm, rigid-body
dynamics $M(q)\ddot q+C(q,\dot q)\dot q+g(q)=u$ at
$\Delta t=0.02$\,s, and end-effector cost
$\ell=\|\mathrm{FK}(q)-p_{\mathrm{goal}}\|_Q^2+\|u\|_R^2$ over
$H\!=\!64$ ($\sim$1.3\,s). Each episode places 2--3 cylindrical
obstacles on the table and draws $p_{\mathrm{goal}}$ uniformly from
a $0.4\!\times\!0.4\!\times\!0.2$\,m box, exposing visibly multimodal
arm postures. Demonstrations are taken from \texttt{Robomimic}'s
\texttt{can}/\texttt{square} proficient-human splits, retargeted to
torque control by inverse dynamics. As an analytical companion we
use a 2D double-integrator $\dot x=Ax+Bu$ with quadratic cost,
where the DARE solution $-P_t x_t$ is closed-form ground truth for
the importance-sampling estimator.

\subsection{Plan for Implementation}

% Describe in \textit{one paragraph} how you will implement its
% simulation and conduct numerical experiments.

Simulation in \texttt{MuJoCo}, learning in PyTorch on a single GPU.
$v_\theta$ is a 1D temporal U-Net ($\sim$5\,M params) over
$(q,\dot q,u)$ trajectories with FiLM conditioning on
$p_{\mathrm{goal}}$, trained with Adam (lr $10^{-4}$, batch 64,
$\sim$10$^5$ steps) on $N\!\sim\!10^4$ \texttt{Robomimic} trajectories;
inference uses $T\!=\!50$ denoising steps and $M\!\in\![64,512]$
batched rollouts per step. Two experiments map the offline--online
tradeoff: a $\lambda$-sweep with endpoints $\lambda\!=\!0,\infty$
recovering the data-only~\cite{janner2022,flowmp} and MBD-only
baselines, and an OOD test with novel goals, obstacles, and $\pm
30\%$ mass/friction perturbations. We report success rate, mean
cost, and collision rate over 100 episodes.

% -------------------------------------------------------
\bibliographystyle{IEEEtran}
\bibliography{references}

\end{document}
